\section{Introduction}
Hardware/software co-design requires a partition of a task graph into software-executed and hardware-accelerated tasks such that end-to-end latency is minimized without violating the available hardware area. Although the partition is often represented as a binary labeling over nodes, the final quality of that labeling is determined by scheduling. Moving a task to hardware may shorten its local execution time, but it may also introduce communication overhead, alter the critical path, or change the serialization pattern on the software processor. This tight coupling between partitioning and scheduling is one of the central reasons the HW/SW partitioning problem remains difficult \cite{arato2005algorithmic,hou2019survey,teich2012hardware}.

Classical exact methods are useful on small graphs, but their scalability is limited when realistic precedence and resource constraints are modeled together \cite{arato2005algorithmic,hou2019survey}. Metaheuristics improve scalability by searching the binary placement space through repeated black-box evaluations, but they still depend on many expensive schedule evaluations and may struggle when the objective landscape is highly nonconvex. Graph-based learning offers a more structural alternative. In particular, Zheng \textit{et al.} introduced a graph-convolution-based HW/SW partitioner that uses graph features and scheduling feedback to improve partition quality \cite{zheng2021hardware}. The current repository follows that direction and extends it toward a differentiable, order-aware formulation.

This paper is centered on \diffgnnorder, the order-aware implementation exposed in the code as \texttt{diff\_gnn\_order}. The method predicts both relaxed hardware-assignment probabilities and software execution priorities. These outputs are fed into a differentiable scheduler surrogate that approximates makespan under precedence, communication, and single-processor software serialization. During inference, the relaxed solution is decoded, repaired to satisfy the area constraint, optionally refined with local search, and evaluated with both learned and static software priorities. In contrast, the placement-only \diffgnn model does not predict an execution order and is therefore treated here as an ablation rather than as the main method.

\noindent\textbf{Our contributions are as follows:}
\begin{itemize}
\item We formulate the proposed method as an order-aware differentiable graph-neural model for HW/SW partitioning, using the same notation and scheduling-oriented objective structure as the ICCAD draft.
\item We document the concrete \diffgnnorder configuration used in the repository, including the feature construction, edge-weight mechanism, ordering relaxation, and post-processing pipeline used during evaluation.
\item We define an evaluation protocol that separates DAG makespan, static-priority makespan, and learned-priority makespan, and we make the reporting policy explicit: static-priority makespan for all non-order-aware methods, and $\min(M_{\mathrm{static}}, M_{\mathrm{learned}})$ for \diffgnnorder.
\item We describe the repository datasets in a paper-ready format, including TOSA-derived real graphs, the cached SqueezeNet-TOSA instance sweep, and a new SqueezeNet-like synthetic generator for 1,000- and 10,000-node stress tests.
\end{itemize}

The remainder of the paper is organized as follows. Section~\ref{sec:preliminaries} formulates the problem and its differentiable relaxation. Section~\ref{sec:related} reviews prior exact, metaheuristic, and graph-based approaches. Section~\ref{sec:method} presents \diffgnnorder and clarifies how \diffgnn is used as an ablation. Section~\ref{sec:settings} details the experimental settings, Section~\ref{sec:results} provides the results templates and reporting structure, and Section~\ref{sec:conclusion} concludes.
